%============================================================
% Chapter 12: Designing Array-Based Data Structures
% STATUS: PLANNED  (generated by tools/outline.py; replace this file when the chapter is written)
%============================================================
\chapter{Designing Array-Based Data Structures}
\label{ch:design}

\begin{tcolorbox}[enhanced, colback=white, colframe=secondaryblue, boxrule=1pt, arc=2mm]
\textbf{\large\textcolor{primaryblue}{Arrays are the engine inside most data structures: dynamic arrays, $\bigO{1}$ deletion by swapping, lazy resets with timestamps, versioned arrays and array-backed hash sets.}}
\end{tcolorbox}

\chapterfacts{1200}{2000}{5}{10}{30}

\begin{plannedbox}[Planned Contents --- this chapter is not written yet]
\textbf{Part A --- Lessons}
\begin{enumerate}[leftmargin=*, label=\textbf{A\arabic*.}, itemsep=1pt]
    \item The dynamic array from scratch: growth, shrinking and the amortised bound
    \item Delete in $\bigO{1}$ by swapping with the last element; index maps for random access
    \item Lazy initialisation and resets in $\bigO{1}$: timestamps and the sparse-set trick
    \item Versioned arrays: snapshots by per-cell history and binary search (full persistence is the tree book)
    \item Open addressing: a hash set in a plain array; bit arrays and \texttt{vector<bool>}
\end{enumerate}
\textbf{Part B --- Patterns} (with candidate problems, checked when the chapter is written)
\begin{enumerate}[leftmargin=*, label=\textbf{P\arabic*.}, itemsep=1pt]
    \item \textbf{Swap-with-last deletion} --- Insert Delete GetRandom O(1) {\small\color{darkgray}(LeetCode 380)}, Insert Delete GetRandom O(1) - Duplicates allowed {\small\color{darkgray}(LeetCode 381)}
    \item \textbf{Snapshots and versions} --- Snapshot Array {\small\color{darkgray}(LeetCode 1146)}
    \item \textbf{Counters of counters} --- Frequency Tracker {\small\color{darkgray}(LeetCode 2671)}
    \item \textbf{Array-backed hashing} --- Design HashSet {\small\color{darkgray}(LeetCode 705)}
    \item \textbf{Lazy resets with timestamps}
\end{enumerate}
\textbf{Connections.} Circular buffers and queues are the stack-queue-heap book; persistent arrays via segment trees are the tree book.
\end{plannedbox}
